\section*{Bab \ref{ch:b3:lithosphere-thermal-models} --- Model Termal Litosfer}

\begin{solution}{exo:b3:lithosphere-thermal-models:1}
$L^2/\kappa = \qty{e12}{s}$, \qty{e14}{s} dan \qty{e16}{s}: 0.032, 3.2 dan
\qty{317}{Ma}. Satu kilometer melupakan suhunya dalam puluhan ribu
tahun, seratus kilometer dalam ratusan juta tahun.
\end{solution}

\begin{solution}{exo:b3:lithosphere-thermal-models:2}
$u = \qty{e-3}{m}/\qty{3.156e7}{s} = \qty{3.17e-11}{m.s^{-1}}$, sehingga
$\mathrm{Pe} = 3.17\times10^{-11}\times\num{3e4}/10^{-6} = 0.95$. Adveksi dan
konduksi sama pentingnya: batuan yang naik membawa panas ke atas
dan membuat geoterm lebih curam di dekat permukaan, sedangkan konduksi saja akan
meremehkan gradiennya.
\end{solution}

\begin{solution}{exo:b3:lithosphere-thermal-models:3}
$\sqrt{\kappa t} = \sqrt{\num{0.8e-6}\times5\times\num{3.156e13}} =
\qty{11236}{m}$, $\eta = 10000/(2\times11236) = 0.445$, dan $T = 1450\times0.471
= \qty{683}{\celsius}$.
\end{solution}

\begin{solution}{exo:b3:lithosphere-thermal-models:4}
$490/\sqrt1 = 490$, $490/5 = 98$ dan $490/10 = \qty{49}{mW.m^{-2}}$. Sonda
paling mungkin meremehkan yang pertama: kerak muda retak dan terbuka bagi
air laut yang bersirkulasi, yang membawa sebagian panas keluar melalui mata air panas.
\end{solution}

\begin{solution}{exo:b3:lithosphere-thermal-models:5}
$z = 2\times1.163\sqrt{\kappa t}$: dengan $\sqrt{\kappa t} = 15.89$, 31.78 dan
\qty{50.25}{km}, $z = 37.0$, 73.9 dan \qty{116.9}{km}. Lapisan yang mendingin
menjadi dua kali lipat dari 10 hingga \qty{40}{Ma}; pada \qty{100}{Ma} ia sudah
lebih dalam daripada lempeng setebal \qty{95}{km}.
\end{solution}

\begin{solution}{exo:b3:lithosphere-thermal-models:6}
$C$ sebanding dengan $\alpha$: $\alpha = \num{3.1e-5}\times350/370.5 =
\qty{2.93e-5}{K^{-1}}$. $C$ juga sebanding dengan $T_m - T_0$ (mantel pada
$1450\times350/370.5 = \qty{1370}{\celsius}$) dan terhadap $\sqrt\kappa$ (suatu
difusivitas $0.8\times(350/370.5)^2 = \qty{0.71e-6}{m^2.s^{-1}}$): kedalaman saja
tidak dapat memisahkannya, itulah sebabnya aliran panas juga digunakan.
\end{solution}

\begin{solution}{exo:b3:lithosphere-thermal-models:7}
$3.1\times1450/\num{95e3} = \qty{47.3}{mW.m^{-2}}$ dan $3.1\times1350/\num{125e3}
= \qty{33.5}{mW.m^{-2}}$. Setengah-ruang mencapainya pada $(490/47.3)^2 =
\qty{107}{Ma}$ dan $(490/33.5)^2 = \qty{214}{Ma}$: lempeng yang tipis dan panas mempertahankan
aliran panas dasar samudra tua sebesar \qty{47}{mW.m^{-2}} sejak sekitar \qty{100}{Ma}
seterusnya, sedangkan lempeng yang lebih tebal dan lebih dingin membuatnya turun di bawah
\qty{40}{mW.m^{-2}}, nilai setengah-ruang pada \qty{150}{Ma}.
\end{solution}

\begin{solution}{exo:b3:lithosphere-thermal-models:8}
$A_0h_r = 0.4\times60 = \qty{24}{mW.m^{-2}}$, sehingga $A_0 =
\num{24e-3}/\num{e4} = \qty{2.4}{\micro W.m^{-3}}$ dan $q_m = 60 - 24 =
\qty{36}{mW.m^{-2}}$. Pada \qty{35}{km}, $T = \num{36e-3}\times\num{35e3}/2.5 +
\num{2.4e-6}\times\num{e8}/2.5\times(1 - e^{-3.5}) = 504 + 96\times0.970 =
\qty{597}{\celsius}$.
\end{solution}

\begin{solution}{exo:b3:lithosphere-thermal-models:9}
$h_r = \Delta q_0/\Delta A_0 = \num{18e-3}/\num{1.8e-6} = \qty{10}{km}$, dan
$q_m = q_0 - h_rA_0 = 66 - 30 = \qty{36}{mW.m^{-2}}$, karena $h_rA_0 =
\num{e4}\times\qty{3.0e-6}{W.m^{-3}} = \qty{30}{mW.m^{-2}}$.
\end{solution}

\begin{solution}{exo:b3:lithosphere-thermal-models:10}
(a) $Q(t) = \int_0^tk(T_m - T_0)/\sqrt{\pi\kappa t'}\,\mathrm dt' =
2k(T_m - T_0)\sqrt{t/(\pi\kappa)}$. (b) Panas yang hilang adalah kandungan panas
yang dikeluarkan dari kolom, $Q = \rho_mc\int_0^\infty(T_m - T)\,\mathrm dz$, dan
isostasi memberikan $(\rho_m - \rho_w)(d - d_r) = \rho_m\alpha\int_0^\infty(T_m -
T)\,\mathrm dz = \alpha Q/c$. (c) $c = 3.1/(3300\times\num{0.8e-6}) =
\qty{1174}{J.kg^{-1}.K^{-1}}$; pada \qty{40}{Ma}, $t = \qty{1.262e15}{s}$ dan
$Q = 2\times3.1\times1450\times\sqrt{\num{1.262e15}/(\pi\times\num{0.8e-6})} =
\qty{2.015e14}{J.m^{-2}}$; $d - d_r = \num{3.1e-5}\times\num{2.015e14}/
(1174\times2270) = \qty{2343}{m}$, yaitu $370.5\sqrt{40}$.
\end{solution}

\begin{solution}{exo:b3:lithosphere-thermal-models:11}
Suku $n = 3$ adalah $3104\times(8/\pi^2)\times\tfrac19e^{-9t/36.2} =
280\,e^{-9t/36.2}$ meter: \qty{218}{m} pada \qty{1}{Ma}, \qty{23}{m} pada
\qty{10}{Ma} dan \qty{1.9}{m} pada \qty{20}{Ma}. Nilai itu turun menjadi \qty{10}{m} ketika
$e^{-9t/36.2} = 10/280$, pada $t = (36.2/9)\ln28 = \qty{13.4}{Ma}$.
\end{solution}

\begin{solution}{exo:b3:lithosphere-thermal-models:12}
(a) Untuk $z \gg h_r$, $e^{-z/h_r} \to 0$. (b) $q_m = \qty{24}{mW.m^{-2}}$,
$A_0h_r^2/k = \num{16e-3}\times\num{e4}/2.5 = \qty{64}{K}$, dan $z = (1300 -
64)\times2.5/\num{24e-3} = \qty{128.8}{km}$. (c) $q_m = 2.5\times1236/\num{250e3}
= \qty{12.4}{mW.m^{-2}}$: \omterm{def:b3:lithosphere-thermal-models:reduced}{aliran panas tereduksi} harus kira-kira dibagi dua, seperti di bawah
kraton yang akar tebal dan dinginnya membiarkan sedikit panas masuk dari bawah.
\end{solution}

\begin{solution}{pb:b3:lithosphere-thermal-models:1}
\textbf{1.} $\sqrt{\num{0.8e-6}\times\num{3.156e14}} = \qty{15889}{m}$ dan
$\sqrt{\num{0.8e-6}\times\num{1.2624e15}} = \qty{31778}{m}$.

\textbf{2.} $\eta = 20000/31778 = 0.6294$ dan $20000/63556 = 0.3147$:
$T = 1450\times0.6266 = \qty{909}{\celsius}$ pada \qty{10}{Ma} dan
$1450\times0.3437 = \qty{498}{\celsius}$ pada \qty{40}{Ma}.

\textbf{3.} $z = 2\times0.7174\sqrt{\kappa t}$: \qty{22.8}{km} dan
\qty{45.6}{km}. Ia mencapai \qty{60}{km} ketika $\sqrt{\kappa t} = 60/1.4348 =
\qty{41.82}{km}$, pada $t = (\num{41.82e3})^2/\num{0.8e-6} = \qty{2.186e15}{s}$,
\qty{69.3}{Ma}.

\textbf{4.} $q = 3.1\times1450/(\sqrt\pi\times15889) = \qty{159.6}{mW.m^{-2}}$ pada
\qty{10}{Ma} dan \qty{79.8}{mW.m^{-2}} pada \qty{40}{Ma}; $C_q = 159.6\sqrt{10}
= \qty{505}{mW.m^{-2}.Ma^{1/2}}$.

\textbf{5.} $u = \qty{7.92e-10}{m.s^{-1}}$, $\mathrm{Pe} = 7.92\times10^{-10}
\times\num{95e3}/\num{0.8e-6} = 94$. Panas bergerak ke samping bersama lempeng
seratus kali lebih cepat daripada ia berdifusi ke samping, sehingga setiap
kolom mendingin hanya melalui konduksi vertikal.

\textbf{6.} $(\num{95e3})^2/\num{0.8e-6} = \qty{1.13e16}{s}$, \qty{357}{Ma}.
Pada \qty{150}{Ma}, $\sqrt{\kappa t} = \qty{61.5}{km}$ dan lapisan yang mendingin,
beberapa kali lipat darinya, mencapai jauh di bawah \qty{95}{km}: model setengah-ruang untuk dasar samudra yang setua itu
mengandaikan bahwa dingin telah mencapai kedalaman tempat mantel berkonveksi.

\textbf{7.} $(\rho_m - \rho_w)(d - d_r) = \rho_m\alpha\int(T_m - T)\,\mathrm dz$,
massa kolom punggungan dan kolom tua sama besar di atas
kedalaman kompensasi. $T_m - T = T_m\operatorname{erfc}\eta$ turun di bawah
\qty{0.5}{\percent} dari $T_m$ untuk $\eta > 2$, sehingga bagian yang dalam tidak menambah apa-apa.

\textbf{8.} $C = 2\times3300\times\num{3.1e-5}\times1450/2270\times
\sqrt{\num{0.8e-6}\times\num{3.156e13}/\pi} = 0.13069\times2834.8 =
\qty{370.5}{m.Ma^{-1/2}}$; $d = 2500 + 370.5\sqrt{10} = \qty{3672}{m}$ dan
$2500 + 370.5\sqrt{40} = \qty{4843}{m}$.

\textbf{9.} Model memberikan $2500 + 370.5\sqrt{42.5} = \qty{4915}{m}$,
\qty{440}{m} lebih dalam daripada median. Median memasukkan beban
sedimen, gunung laut dan dataran vulkanik, tonjolan titik panas, serta
sel-sel di lereng zona rekahan dan tepi, yang semuanya membuat dasar
lebih dangkal.

\textbf{10.} $C$ sebanding dengan $T_m$: $370.5\times1350/1450 =
\qty{344.9}{m.Ma^{-1/2}}$, mendekati 350 hasil taksiran.

\textbf{11.} Air laut yang bersirkulasi melalui kerak muda yang retak membawa
sebagian panas keluar lewat mata air panas, di antara sensor.

\textbf{12.} $\tau = (\num{95e3})^2/(\pi^2\times\num{0.8e-6}) =
\qty{1.143e15}{s} = \qty{36.2}{Ma}$; amplitudonya adalah $3300\times
\num{3.1e-5}\times1450\times\num{95e3}/4540 = \qty{3104}{m}$.

\textbf{13.} $e^{-40/36.2} = 0.33124$, $e^{-360/36.2}/9 = \num{5.3e-6}$ dan
suku $n = 5$ jauh lebih kecil lagi; kurung bernilai $1 - 0.81057\times0.33125 =
0.73150$, dan $d = 2500 + 3104\times0.73150 = \qty{4771}{m}$.

\textbf{14.} Dari suku $n = 3$, $3104\times0.81057\times\num{5.3e-6} =
\qty{0.013}{m}$: pada \qty{40}{Ma}, $e^{-9t/\tau}$ bernilai $10^{-4}$ dari
$e^{-t/\tau}$, dan hanya mode pertama yang tersisa.

\textbf{15.} Pada \qty{40}{Ma}, $4843 - 4771 = \qty{73}{m}$; pada
\qty{50}{Ma}, $2500 + 370.5\sqrt{50} = 5120$ dibandingkan $2500 + 3104\times(1 -
0.81057\times0.25127) = 4972$: \qty{148}{m}. Dengan interpolasi, $40 + 10\times
(100 - 73)/(148 - 73) = \qty{43.6}{Ma}$, sekitar \qty{44}{Ma} (deret penuh
memberikan 44.0).

\textbf{16.} $q = (3.1\times1450/\num{95e3})[1 + 2(0.33124 + 0.01204 +
0.00005)] = 47.32\times1.68666 = \qty{79.8}{mW.m^{-2}}$, sama seperti
setengah-ruang: pada \qty{40}{Ma} dasar belum terasa di
permukaan.

\textbf{17.} $2500 + 3104 = \qty{5604}{m}$ dan $kT_m/L = \qty{47.3}{mW.m^{-2}}$.

\textbf{18.} $d(120) = 2500 + aL\bigl[1 - (8/\pi^2)e^{-120/\tau}\bigr]$, dengan
$a = \rho_m\alpha T_m/\bigl(2(\rho_m - \rho_w)\bigr) = \qty{32.673}{m/km}$ dan
$\tau = L^2/(\pi^2\kappa)$.

\textbf{19.} $L = 3200/\bigl(32.673\,[1 - 0.81057e^{-120/\tau}]\bigr)$. Dari
$L = \qty{100}{km}$: $\tau = \qty{40.13}{Ma}$, kurung 0.95925, $L =
\qty{102.10}{km}$; lalu $\tau = 41.83$, $L = 102.67$; $\tau = 42.30$, $L =
102.82$; $\tau = 42.43$, $L = 102.87$; nilainya menetap di $L = \qty{102.9}{km}$.

\textbf{20.} $\tau = \qty{42.48}{Ma}$, $aL = \qty{3362}{m}$: $d(40) = 2500 +
3362\times(1 - 0.81057e^{-40/42.48}) = \qty{4799}{m}$, dibandingkan \qty{4843}{m}
untuk setengah-ruang dan \qty{4800}{m} terukur: lempeng ini cocok dengan kedua umur.

\textbf{21.} $\tau = \qty{42.5}{Ma}$; kedalaman asimtotik $2500 + 3362 =
\qty{5862}{m}$.

\textbf{22.} $3362\times0.81057\,e^{-t/42.48} = 100$ memberikan $t = 42.48\ln27.25 =
\qty{140}{Ma}$.

\textbf{23.} Lempeng yang cocok dengan dasar tua samudra fiktif
\textbf{\qty[text-series-to-math=true]{102.9}{km}} tebalnya.
\end{solution}
